\mysec{Lista 4}
\etocsetnexttocdepth{subsection}
\localtableofcontents
\newcounter{quatro}
\stepcounter{quatro}

%ex 1
\myexe{\thequatro}
\bigline
\textbf{Seja $(X, \mcal F, \mu)$ um espaço de medida e $f \in L^p(X, \mu)$, 
com $1 \le p < +\infty$. Prove a Desigualdade de Chebyshev:
\begin{align*}
  \mu\parth{\set{x \in X \st |f(x)| > \alpha}} \le \cparth{\frac{\pnorm p f}{\alpha}}^p, \quad \forall \alpha>0.
\end{align*}}
\begin{proof}[\textbf{Dem:}] Como $f \in L^p (X, \mu)$, temos $|f|^p \in L^1(X, \mu)$, ou seja, 
  \begin{align*}
    \int |f|^p d\mu < + \infty.
  \end{align*}
  Mais ainda, como $|f|^p \ge 0$, temos que em particular que todos os representantes da classe 
  $[|f|^p]$ são mensuráveis não-negativos. Portanto, 
  da desigualdade de Chebyshev para funções em $\fplus M (X, \mcal F)$, dados $\beta = \alpha^p > 0$ e
  $E_\beta = \set{x \in X \st |f(x)|^p > \beta}$,
  \begin{align*}
    \mu(\set{x \in X \st |f(x)| > \alpha}) = \mu(E_\beta) \le \frac{1}{\beta} \int |f|^p d\mu 
    = \frac{1}{\alpha^p} \pnorm p f^p.
  \end{align*}
\end{proof}
\bigline

%ex 2 
\stepcounter{quatro}
\myexe{\thequatro}

\noindent 
\textbf{Sejam $(X, \mcal F, \mu)$ um espaço de medida e $f \in L^p (X, \mu)$, com $1 \le p < +\infty$. 
Mostre que o conjunto $E = \set{x \in X \st |f(x)| \neq 0}$ é $\sigma$-finito.}
\begin{proof}[\textbf{Dem:}] Note que $|f(x)|^p \ge 0$ para cada 
  $x \in X$ e $1 \le p \le +\infty$. Para $p \neq +\infty$, considere  
  $E_n = \set{x \in X \st |f(x)| > \frac{1}{n}}$. 
  Note que $E = \bigcup_{n \in \bN} E_n$, desse modo, vejamos que 
  $E_n$ é finito. Note que da desigualde de Chebyshev, para 
  cada $n \in \bN$,
  \begin{align*}
    \mu(E_n) \le \cparth{\frac{||f||_p}{n}}^p < +\infty.
  \end{align*}
  Com isso, segue que $E$ é $\sigma$-finito.
\end{proof}

%ex 3
\nextexe{quatro}
\textbf{Sejam $(X, \mcal F, \mu)$ um espaço de medida e $f \in L^p(X, \mu)$, 
com $1 \le p \le +\infty$. Definindo $E_n = \set{x \in X \st |f(x)| \ge n}$, 
mostre que $\lim \mu(E_n) = 0$.}
\begin{proof}[\textbf{Dem:}] Note que, para cada $n \in \bN$ temos 
$E_n \supseteq E_{n+1}$. Mais ainda, se $p \neq +\infty$, para cada $n \in \bN$, 
  da desigualdade de Chebyshev para espaços $L^p(X, \mu)$, dado $f \in L^p(X, \mu)$, 
\begin{align*}
  \mu(E_n) \le \cparth{\frac{||f||_p}{n}}^p \implies 
  \lim_{n \rightarrow \infty} \mu(E_n) \le 
  \lim_{n \rightarrow \infty} \cparth{\frac{||f||_p}{n}}^p = 0. 
\end{align*}
Se $p = +\infty$, por definição, dado $f \in L^\infty(X, \mu)$, 
\begin{align*}
  ||f||_\infty = \inf \set{a \ge 0 \st \mu(\set{x \in X \st |f(x)| > a}) = 0},
\end{align*}
  isto é, $\mu(\set{x \in X \st |f(x)| > ||f||_\infty}) = 0$.  
  Visto que $f \in L^\infty(X, \mu)$, existe $M \in \bN$ tal que 
$||f||_\infty < M$. Desse modo, para todo $n \ge M$,
\begin{align*}
  \mu(\set{x \in X \st |f(x)| \ge n}) \le 
  \mu\parth{\set{x \in X \st |f(x)| > ||f||_\infty}} = 0,
\end{align*}
  portanto, $\lim \mu(E_n) = 0$.
\end{proof}

%ex 4
\newexe{quatro}
\textbf{Sejam $X = \bN$ e $\mu$ a medida da contagem em $\bN$. 
Considere $f: \bN \rightarrow \bR$ definida por $f(n)=1/n$. Mostre que 
$f$ não pertence a $L^1(\bN, \mu)$, mas pertence a $L^p(\bN, \mu)$ 
para todo $1 < p \le \infty$.}
\begin{proof}[\textbf{Dem:}] Ora, como visto em exercícios anteriores, 
\begin{align*}
  \int |f| d\mu = \sum_{k \in \bN} |f(k)| = \sum_{k \in \bN} \mparth {\frac{1}{k}} 
  = \sum_{k \in \bN} \frac{1}{k},
\end{align*}
com isso, segue que $\int |f| d\mu$ é igual a série harmônica, i.e., $\int |f| d\mu = +\infty$,
  ou seja, $f \not \in L^1(X, \mu)$. 
  Entretanto, 
\begin{align*}
  \int |f|^p d\mu = \sum_{k \in \bN} |f(k)|^p = \sum_{k \in \bN} \mparth{\frac{1}{k}}^p 
  = \sum_{k \in \bN} \frac{1}{k}^p,
\end{align*}
ou seja, da convergência das $p$-séries, segue que $\int |f|^p d\mu < +\infty$, 
  isto é, $f \in L^p(X, \mu)$ para 
  $1 < p < +\infty$. Agora, para $p = +\infty$, note que, para cada $n \in \bN$, 
\begin{align*}
  f(n) \le \max f(k) = 1 \implies ||f||_\infty \le \max f(k) = 1 < +\infty. 
\end{align*}
  Portanto $f \in L^\infty(X, \mu)$.
\end{proof}

%ex 5
\nextexe{quatro}
\textbf{Sejam $X = \bN$ e $\mu$ a medida da contagem em $\bN$. 
Prove que se $f \in L^p(\bN, \mu)$, então $f \in L^q(\bN, \mu)$ para todo $1 \le p \le q < \infty$, 
e $||f||_q \le ||f||_p$.}
\begin{proof}[\textbf{Dem:}] Note que para $p = q$ segue diretamente o resultado. 
  Desse modo, suponhamos que $f \in L^p(\bN, \mu)$. Já que $\mu$ é a medida da contagem 
  em $\bN$, 
  \begin{align*}
    ||f||_p = \parth{\int |f|^p d\mu}^{1/p} = \parth{\sum_{k \in \bN} |f(k)|^p}^{1/p}
    \ge |f(n)|, \quad \forall n \in \bN.
  \end{align*}
  Daqui, note que, para cada $n \in \bN$, dado $q-p > 0$, 
  \begin{align*}
    |f(n)|^q = |f(n)|^p\cdot |f(n)|^{q-p} \le |f(n)|^p \cdot ||f||_p^{q-p}
  \end{align*}
  Ou seja, 
  \begin{align*}
    \parth{\int |f|^q d\mu}^{1/q} \le \parth{\int |f|^p\cdot ||f||_p^{q-p} d\mu}^{1/q}
    = \parth{||f||_p^{q-p}\int |f|^pd\mu}^{1/q} = ||f||_p,
  \end{align*}
  isto é, $||f||_q \le ||f||_p$. Mais ainda, segue diretamente da desigualdade 
  que $f \in L^q(\bN, \mu)$.

\end{proof}

%ex 6
\nextexe{quatro}
\textbf{Seja $X = \bN$ e defina $\lambda: \mcal P(\mbb N) \rightarrow \bN$ por 
\begin{align*}
  \lambda(E) = \sum_{n \in E}\frac{1}{n^2}
\end{align*}
Considere $f: \bN \rightarrow \bR$ dada por $f(n)= \sqrt{n}$.
\begin{enumerate}[label=\alph*)]
  \item Mostre que $\lambda$ é uma medida em $\bN$ e que $\lambda(\bN) < +\infty$.
    \begin{proof}[\textbf{Dem:}]
      Por definição segue que $\lambda(E) \ge 0$ para todo $E \in \mcal P(\bN)$ não 
      vazio. Para $E = \emptyset$,defina $\lambda(E) = 0$. Daqui, vejamos que vale 
      a $\sigma$-aditividade. Dado $\set{E_j}$ uma sequência em $\mcal F$ dois a 
      dois disjunto, note que, 
      \begin{align*}
        \lambda \parth{\bigcup_{j \in \bN} E_j} = \sum_{n \in \bigcup E_j} \frac{1}{n^2}
        = \sum_{j \in \bN} \cparth{\sum_{n \in E_j} \frac{1}{n^2}} 
        = \sum_{j \in \bN} \lambda(E_j), 
      \end{align*}
      ou seja, $\lambda$ é uma medida em $\bN$. Como $\lambda(\bN)$ é a $p$-série 
      com $p=2$, segue que $\lambda(\bN) < +\infty$.
    \end{proof}
  \item Mostre que $f \in L^p(\bN, \lambda)$ se, e somente se, $1\le p < 2$.
  \begin{proof}[\textbf{Dem:}] Seja $p \ge 1$. Dado qualquer 
    $f: \bN \rightarrow \bR$ segue que para todo $E \subset \bR$, 
    $f^{-1}(E) \in \mathcal P (\bN)$, isto é, toda função $f: \bN \rightarrow \bR$
    é mensurável em $(X, \mathcal P(\bN))$. Daqui, note que dos exercícios anteriores,
    dado $|f|^p \in M(\bN, \mathcal P(\bN))$,
    \begin{align*}
      \int |f|^p d\lambda &= \sum_{n \in \bN} |f(n)|^p \lambda(\set{n})
      = \sum_{n \in \bN}  \cparth{|f(n)|^p \sum_{k \in \set{n}} \frac{1}{k^2}}\\
      &= \sum_{n \in \bN} \frac{|f(n)|^p}{n^2} = \sum_{n \in \bN} \frac{n^{p/2}}{n^2}
      = \sum_{n \in \bN} \frac{1}{n^{2-p/2}}.
    \end{align*}
    Da convergência das p-séries, segue que $\sum_{n \in \bN} \frac{1}{n^p}$ converge
    para algum número real se, 
    e somente se $p>1$. Desse modo, $\int |f|^p d\lambda < + \infty$ se, e somente se 
    $2-p/2 > 1$, isto é, $p < 2$. Portanto, $f \in L^P(\bN, \lambda)$ se, e somente se 
    $1\le p < 2$.
  \end{proof}
\end{enumerate}}

%ex 7
\nextexe{quatro}
\textbf{Seja $(X, \mcal F, \mu)$ um espaço de medida finita. Se $f$ é mensurável em $X$, 
defina $E_n = \set{x \in X \st (n - 1) \le |f(x)| < n}$. Mostre que $f \in L^1(X, \mu)$ 
se, e somente se, 
\begin{align*}
  \sum_{n \in \bN} n \mu(E_n) < +\infty.
\end{align*}
Generalize para $1 \le p < +\infty$, isto é, mostre que $f \in L^p(X, \mu)$, se 
e somente se 
\begin{align*}
  \sum_{n \in \bN} n^p \mu(E_n) < + \infty.
\end{align*}}
\begin{proof}[\textbf{Dem:}] Definindo $N = \{x \in X \st |f(x)| = +\infty \}$, segue que 
  $N^C = \{x \in X \st |f(x)| < +\infty\}$. E, como $N^C = \bigcup E_n$ onde para 
  cada $k, l \in \bN$, $E_k \cap E_l = \emptyset$ sempre que $k \neq l$, temos que 
  $\mu\parth{N^C} = \mu\parth{\bigcup E_n} = \sum \mu(E_n)$. Com isso, vejamos a 
  ida e a volta da proposição.\\
$(\implies)$ Suponhamos que $f \in L^1(X, \mu)$. Então, $|f| < +\infty$ $\mu$-q.t.p.. Mais 
  precisamente $\mu(N) = 0$. Daqui, note que,
\begin{align*}
  \int |f| d\mu &= \int_N |f| d\mu + \int_{N^C} |f| d\mu = \int_{N^C} |f| d\mu 
  = \int |f| \chi_{N^C} d\mu\\
  &= \int |f| \chi_{\bigcup E_n} d\mu = \int |f| \parth{\sum_{n \in \bN} \chi_{E_n}}d\mu < +\infty.
\end{align*}
Ou seja, 
\begin{align*}
  \int |f| d\mu = \int |f| \parth{\sum_{n \in \bN} \chi_{E_n}} = \sum_{n \in \bN} \int |f| \chi_{E_n} d\mu.
\end{align*}
Pela definição de $E_n$, temos,
\begin{align*}
  +\infty > \int |f| d\mu \ge \sum_{n \in \bN} \int (n-1) \chi_{E_n} d\mu
  = \sum_{n \in \bN} (n-1) \mu\parth{E_n}.
\end{align*}
  Como $N^C \subset X$ e $\mu(X) < +\infty$,
  \begin{align*}
    +\infty > \cparth{\sum_{n \in \bN} (n - 1) \mu\parth{E_n}} + 
    \cparth{\sum_{n \in \bN} \mu\parth{E_n}} 
    = \sum_{n \in \bN} n\mu\parth{E_n}.
  \end{align*}
  $(\impliedby)$ Como $f: X \rightarrow \bR$, note que $\mu({x \in X \st |f(x)| = +\infty}) = 0$. Portanto, 
  temos da definição de $f$ que $X = \bigcup E_n$. Já que $f$ é mensurável e limitada por $n$ em $E_n$
  para todo $n \in \bN$,
  \begin{align*}
    \int |f| d\mu = \int |f| \chi_X d\mu = \int |f| \chi_{\bigcup E_n} d\mu = \int \sum_{n \in \bN} |f| \chi_{E_n} d\mu 
    \le \int \sum_{n \in \bN} n \chi_{E_n} d\mu.
  \end{align*}
  Já que para cada $n \in \bN$, $n \chi_{E_n} \in \mathcal M^+(X, \mu)$,
  \begin{align*}
    \int |f| d\mu \le \int \sum_{n \in \bN} n \chi_{E_n} d\mu = \sum_{n \in \bN} \int n \chi_{E_n}d \mu, 
    = \sum_{n \in \bN} n \mu(E_n) < +\infty
  \end{align*}
  ou seja, $f \in L^1(X, \mu)$.

  Agora, vejamos que vale a generalização. \\
  $(\implies)$ Suponhamos $f \in L^p(X, \mu)$ e vejamos que 
  $\sum\limits_{n \in \bN} n^p \mu(E_n) < +\infty$. Como $f \in L^p(X, \mu)$, 
  $\mparth{f}^p < +\infty$ $\mu$-q.t.p., ou seja, $\mu(N) = 0$. Portanto,
  \begin{align*}
    \int \mparth{f} d\mu &= \int_{N} \mparth{f}^p d\mu + \int_{N^C} \mparth{f}^p d\mu 
    = \int_{N^C} \mparth{f}^p d\mu\\ 
    &= \int \mparth{f}^p \parth{\sum_{n \in \bN} \chi_{E_n}} d\mu < +\infty.
  \end{align*}
  Dito isso, 
  \begin{align*}
    +\infty &>\int \mparth{f}^p d\mu = 
    \int \mparth{f}^p \parth{\sum_{n \in \bN} \chi_{E_n}} d\mu 
    = \sum_{n \in \bN} \int \mparth{f}^p\chi_{E_n} d\mu\\ 
    &\ge \sum_{n \in \bN} \int (n-1)^p \chi_{E_n}d\mu = \sum_{n \in \bN} (n-1)^p \mu(E_n).
  \end{align*}
  Sendo assim, note que para $n \ge 2$, temos $n \le 2(n-1)$, ou seja, $n^p \le 2^p(n-1)^p$. 
  Daqui, já que o espaço é finito, $E_1 \subseteq X$ tem medida finita, e, portanto,
  \begin{align*}
    \sum_{n \in \bN} n^p\mu(E_n) &= \mu(E_1) + \sum_{n = 2}^\infty n^p\mu(E_n)
    \le \mu(E_1) + \sum_{n = 2}^\infty 2^p(n-1)^p \mu(E_n)\\ 
    &= \mu(E_1) + 2^p \sum_{n = 2}^\infty (n-1)^p \mu(E_n) < +\infty.
  \end{align*}
  $(\impliedby)$ Suponhamos que $\sum\limits_{n \in \bN} n^p \mu(E_n) < +\infty$. 
  Por hipótese, $f: X \rightarrow \bR$, ou seja, $\bigcup E_n = X$. Com isso, 
  como $f$ é mensurável, temos $\mparth{f}$ mensurável e, já que $\phi(x) = x^p$ é 
  contínua para cada $p \in [1, +\infty)$, temos $\mparth{f}^p \in \mcal M^+(X, \mu)$. 
  Sendo assim, 
  \begin{align*}
    \int \mparth{f}^p d\mu &= \int \mparth{f}^p \chi_{X} d\mu 
    = \int \mparth{f}^p \parth{\sum_{n \in \bN} \chi_{E_n}} d\mu 
    = \sum_{n \in \bN} \int \mparth{f}^p \chi_{E_n} d\mu \\
    &= \sum_{n \in \bN} \int n^p \chi_{E_n} d\mu = \sum_{n \in \bN} n^p \mu(E_n)<+\infty. 
  \end{align*}
  Ou seja, $f \in L^p(X, \mu)$.
\end{proof}

%ex 8
\nextexe{quatro}
\textbf{Sejam $(X, \mcal F, \mu)$ um espaço de medida finita e $f \in L^p(X, \mcal F, \mu)$, $1 \le p < \infty$, 
e $\epsilon > 0$. Mostre que existe um conjunto $E_\epsilon \in \mcal F$ com $\mu(E_\epsilon) < +\infty$ 
tal que se $F \in \mcal F$ e $F \cap E_\epsilon = \emptyset$, então $\pnorm p {f \chi_F} < \epsilon$.}
\begin{proof}[\textbf{Dem:}] Ora, como $(X, \mcal F, \mu)$ é um espaço de medida finita, então, 
  para todo $A \in \mcal F$, segue que $\mu(A) < +\infty$. Agora, como $f \in L^p(X, \mcal F, \mu)$, 
  todo representante de $|f|^p \in L^1(X, \mu)$ é mensurável. Portanto, temos que 
  \begin{align*}
    E_\epsilon = \set{x \in X \st |f(x)|^p \ge \frac{\epsilon^p}{\mu(X)}} \in \mcal F.
  \end{align*}
  Agora, se $F \in \mcal F$ e $F \cap E_\epsilon = \emptyset$, segue que, 
  \begin{align*}
    F \subset \set{x \in X \st |f(x)|^p < \frac{\epsilon^p}{\mu(X)}}.
  \end{align*}
  E, como os representantes de $|f|^p$ são 
  mensuráveis, segue que os representantes de $|f|^p\chi_F = |f \chi_F|^p$ também são. Portanto, 
  \begin{align*}
    \int |f \chi_F|^p d\mu = \int | f |^p \chi_F d\mu < \int \frac{\epsilon^p}{\mu(X)} \chi_F d\mu 
    = \epsilon^p \frac{\int_F d\mu}{\mu(X)} = \epsilon^p \cparth{\frac{\mu(F)}{\mu(X)}} \le \epsilon^p.
  \end{align*}
Ou seja, 
\begin{align*}
  \pnorm p {f\chi_F} = \cparth{\int |f\chi_F|^p d\mu}^{1/p} < \epsilon.
\end{align*}
\end{proof}

%ex 9
\nextexe{quatro}
\textbf{Sejam $(X, \mcal F, \mu)$ um espaço de medida e $f \in L^\infty(X, \mu)$. 
Mostre que se $A < \pnorm \infty f$, então existe um conjunto 
$E \in \mcal F$ com $\mu(E) > 0$ tal que $|f(x)| > A$ para todo $x \in E$.}
\begin{proof}[\textbf{Dem:}] Por definição, 
  \begin{align*}
    \pnorm \infty f = \inf \set{\alpha \ge 0 \st \mu(\set{x \in X \st |f(x)| > \alpha}) = 0}
  \end{align*}  
  Desse modo, como $\alpha = ||f||_\infty$ é o menor real extendido satisfazendo 
  \begin{align*}
    \mu(\set{x \in X \st |f(x)| > \alpha}) = 0,
  \end{align*}
  já que $\pnorm \infty f > A$, então $\mu(\set{x \in X \st |f(x)| > A}) > 0$. Definindo 
  $E = \set{x \in X \st |f(x)| > A}$, segue que para todo $x \in E$, temos $|f(x)| > A$,
  mais ainda, o conjunto satisfaz $\mu(E) > 0$, como queríamos.
\end{proof}

%ex 10
\nextexe{quatro}
\textbf{Sejam $(X, \mcal F, \mu)$ um espaço de medida. Mostre que se $f \in L^P(X, \mu)$, $1 \le p \le \infty$, 
e $g \in L^\infty(X, \mu)$, então $fg \in L^p(X, \mu)$ e $\pnorm p {fg} \le \pnorm p f \pnorm \infty g$.}
\begin{proof}[\textbf{Dem:}] Para $p = \infty$, segue diretamente que, 
  \begin{align*}
    \pnorm \infty {fg} = \inf \set{\alpha \ge 0 \st \mu(\set{x \in X \st |f(x)g(x)| > \alpha}) = 0}.
  \end{align*}
  Entretanto, como 
  \begin{align*}
    \mu(\set{x \in X \st |f(x)| > \pnorm \infty f}) = 0, \quad \mu(\set{x \in X \st |g(x)| > \pnorm \infty g}) = 0, 
  \end{align*}
  então, 
  \begin{align*}
    \mu(\set{x \in X \st |f(x)| |g(x)| > \pnorm \infty f \pnorm \infty g}) 
    \le \mu(\set{x \in X \st |f(x)| > \pnorm \infty f}) = 0.
  \end{align*}
  Ou seja, da definição de infimo, segue que, 
  \begin{align*}
    \pnorm \infty {fg} \le \pnorm \infty f \pnorm \infty g.
  \end{align*}
  Já que $f, g \in L^\infty(X, \mu)$, temos da desigualdade acima que $fg \in L^\infty(X, \mu)$. \\ 
  Suponhamos agora que $ 1 \le p < \infty$, Note que, como os representantes 
  de $|f|^p \in L^p(X, \mu)$ são mensuráveis, segue que para todo $\alpha \in \bR$
  temos que $\alpha |f|^p$ é também mensurável.
  Portanto, como $g \in L^\infty(X, \mu)$, sabemos que $g$ é mensurável e 
  por conseguinte, que $|g|$ também é. Já que $\phi: \bR \rightarrow \bR$ definida 
  por $\phi(x) = x^p$ é, em particular para $1 \le p < \infty$, contínua, 
  segue que $|g|^p$ é mensurável, (ver exercício 10 Lista 1). Portanto, 
  $|fg|^p$ é mensurável, desse modo, já que $|fg|^p \le |f|^p\pnorm \infty g ^p$ $\mu$-q.t.p.,
  \begin{align*}
    \int |fg|^P d\mu &= \int |f|^p |g|^p d\mu \le \int |f|^p \pnorm \infty g^p d\mu \\ 
    &= \pnorm \infty g^p \int |f|^p d\mu,
  \end{align*}
  já que $|f|^p \in L^1(X, \mu)$, segue que $\int |fg|^p d\mu \le \pnorm \infty g^p 
  \pnorm p f^p < +\infty$, isto é, $|fg| \in L^p(X, \mu)$. Mais ainda, 
  como $\pnorm p {fg}^p = \int |fg|^p d\mu \le \cparth{\pnorm p f \pnorm \infty g}^p$,
  temos 
  $$\pnorm p {fg} \le \pnorm p f \pnorm \infty g.$$ 
  Como queríamos.
\end{proof}

\nextexe{quatro}
\textbf{Sejam $(X, \mcal F, \mu)$ um espaço de medida e $1 \le p < r \le +\infty$. 
Defina em $L^p(X, \mu) \cap L^r(X, \mu)$ a função $\pnorm {} {\,\cdot\,}:
L^p(X, \mu) \cap L^r(X, \mu) \rightarrow \bR$ dada por 
$\pnorm {} f \coloneq \pnorm p f + \pnorm r f$. Mostre que 
$\pnorm {} {\, \cdot \,}$ é uma norma e que $L^p(X,\mu)\cap L^r(X, \mu)$ é 
um Espaço de Banach com esta norma.}
\begin{proof}[\textbf{Dem:}]
Vejamos primeiro que $\pnorm {} {\, \cdot \,}$ é uma norma em $L^p(X, \mu) \cap L^r(X, \mu)$.
\begin{itemize}
  \item $\snorm f \ge 0$ e $\snorm{f} = 0 \iff f = 0$:\\
    Dado $f \in L^p(X, \mu) \cap L^r(X,\mu)$, como $\pnorm p {\, \cdot \,}$ e 
    $\pnorm r {\, \cdot \,}$ em $L^p(X, \mu) \cap L^r(X, \mu)$, segue que 
    $\snorm f = \pnorm p f + \pnorm r f \ge 0$. Agora, suponhamos que $\snorm f = 0$. Então, 
    \begin{align*}
      \snorm f = \pnorm p f + \pnorm r f = 0 \implies \pnorm p f = - \pnorm r f \implies 
      \pnorm p f = 0 \implies f = 0.
    \end{align*}
    Por outro lado, suponhamos que $f = 0$, desse modo, 
    \begin{align*}
      f = 0 \implies \pnorm p f = \pnorm r f = 0 \implies 
      \snorm f = \pnorm p f + \pnorm r f =0.
    \end{align*}
  \item Dados $\lambda \in \bR$ e $f \in L^p(X, \mu)\cap L^r(X, \mu)$,
    $\snorm{\lambda f} = \mparth{\lambda} \snorm f$:

    Ora, como $\pnorm p {\, \cdot \,}$ e 
    $\pnorm r {\, \cdot \,}$ são normas em $L^p(X, \mu) \cap L^r(X, \mu)$, 
    \begin{align*}
      \mparth{\mparth{\lambda f}} = \pnorm p {\lambda f} + \pnorm r {\lambda f}
      = \mparth \lambda \pnorm p f + \mparth \lambda \pnorm r f = \mparth \lambda \snorm f.
    \end{align*}
  \item Dados $f, g \in L^p(X,\mu)\cap L^r(X, \mu)$, $\snorm {f + g} \le \snorm f + \snorm g$:
    
    Novamente, como $\pnorm p {\, \cdot \,}$ e 
    $\pnorm r {\, \cdot \,}$ são normas em $L^p(X, \mu) \cap L^r(X, \mu)$, 
    \begin{align*}
      \snorm {f+g} &= \pnorm p {f + g} + \pnorm r {f + g} \le 
      \pnorm p f + \pnorm p g + \pnorm r f + \pnorm r g \\ 
      &= \parth{\pnorm p f + \pnorm r f} + \parth{\pnorm p g + \pnorm r g}
      = \snorm f + \snorm g. 
    \end{align*}
\end{itemize}
  Com isso, segue que $\snorm {\, \cdot \,}$ é norma em $L^p(X, \mu) \cap L^r(X, \mu)$.
  Daqui, vejamos que é também um Espaço de Banach quando munido da norma $\snorm {\, 
  \cdot \,}$. Para tal, vejamos que é completo, isto é, toda sequência 
  de Cauchy em $(L^p(X, \mu) \cap L^r(X, \mu), \snorm {\, \cdot \,})$ converge 
  no sentido $\snorm {\, \cdot \,}$ para alguém em $L^p(X, \mu) \cap L^r(X, \mu)$.
  Considere $(f_n)_{n \in \bN}$ uma sequência de Cauchy em $(L^p(X, \mu) \cap L^r(X, \mu), \snorm {\, \cdot \,})$, 
  ou seja, para todo $\epsilon > 0$, existe $N \in \bN$ tal que dados inteiros $n, m \ge N$, temos,
  \begin{align*}
    \snorm{f_n - f_m} < \epsilon.
  \end{align*}
  Ora, por definição, 
  \begin{align*}
    \pnorm p {f_n - f_m} &\le \pnorm p {f_n - f_m} + \pnorm r {f_n - f_m}  = \snorm {f_n - f_m} < \epsilon, \text{ e também, } \\
    \pnorm r {f_n - f_m} &\le \pnorm p {f_n - f_m}  + \pnorm r {f_n - f_m} = \snorm {f_n - f_m} < \epsilon.
  \end{align*}
  ou seja, $(f_n)_{n \in \bN}$ é Cauchy tanto em $(L^p(X, \mu), \pnorm p {\, \cdot \,})$, quanto 
  em $(L^r(X, \mu), \pnorm r {\, \cdot \,})$. Como ambos são Espaços de Banach, temos que, 
  $(f_n)_{n \in \bN}$ converge para $g \in L^p(X, \mu)$ e $h \in L^r(X, \mu)$  
  simultaneamente. Por modos de convergência,
  sabemos que, $(f_n)$ converge em medida 
  para $g$ e $h$ simultaneamente. Das propriedades de convergência em medida, 
  segue que $g = h$ $\mu$-q.t.p.. Desse modo, $g = h = f \in L^p(X, \mu)\cap L^r(X, \mu)$.
  Daqui, note que, para todo $\epsilon > 0$ existe $N \in \bN$ tal que 
  $\pnorm p {f_n - f} < \epsilon/2$ e $\pnorm r {f_n - f}< \epsilon /2 $, ou seja,
  \begin{align*}
    \snorm {f_n - f} = \pnorm p {f_n - f} + \pnorm r {f_n - f} < \epsilon,
  \end{align*}
  isto é, $(f_n)$ converge no sentido $\snorm {\, \cdot \,}$ para 
  $f \in L^p(X, \mu) \cap L^r(X, \mu)$, i.e., 
  $(L^p(X, \mu)\cap L^r(X, \mu), \snorm {\, \cdot \,})$ 
  é um Espaço Completo.
\end{proof}

%ex 12 
\nextexe{quatro}
\textbf{Sejam $(X, \mcal F, \mu)$ um espaço de medida e $1 \le p < r \le +\infty$. 
Defina em $L^p(X, \mu) + L^r(X, \mu)$ a função 
$\snorm {\, \cdot \,}: L^p(X, \mu) + L^r(X, \mu) \rightarrow \bR$ dada por 
\begin{align*}
  \snorm f \coloneq \inf \set{\pnorm p g + \pnorm r h \st 
  g \in L^p(X, \mu), h \in L^r (X, \mu), \text{ e } f = g+ h}
\end{align*}
Mostre que $\snorm {\, \cdot \,}$ é uma norma e 
que $L^p(X, \mu) + L^r(X, \mu)$ é um Espaço de Banach com esta norma.}
\begin{proof}[\textbf{Dem:}] Considere 
\begin{align*}
  A_f = \set{\pnorm p g + \pnorm r h \st g \in L^p(X, \mu), h \in L^r(X, \mu) 
  \text{ e } f = g + h}.
\end{align*}
  Vejamos que $\snorm{\, \cdot \,}$ 
  é norma em $L^p(X, \mu) + L^r(X, \mu)$. 
\begin{itemize}
  \item $\snorm f \ge 0$ e, $\snorm f = 0$ se, e somente se $f = 0$.\\ 
    Ora, note que para todo $g \in L^p(X,\mu)$ e $h \in L^r(X,\mu)$,
    $\pnorm p g \ge 0$ e $\pnorm r h \ge 0$, desse modo, 
    $\pnorm p g + \pnorm r h \ge 0$, em particular, para $F \in A_f$, segue que 
    $F \ge 0$, portanto $\snorm f = \inf A_f\ge 0$. Daqui, suponhamos $\snorm f = 0$
    e vejamos que $f = 0$. Para tal, note que, da definição de infimo, temos que
    dado $\epsilon > 0$,
    existem $g_\epsilon \in L^p(X, \mu)$ e $h_\epsilon \in L^r(X, \mu)$ tais que 
    $\pnorm p {g_\epsilon} + \pnorm r {h_\epsilon} \in A_f$ e satisfazem 
    \begin{align*}
      \pnorm p {g_\epsilon} + \pnorm r {h_\epsilon} < \snorm f + \epsilon.
    \end{align*}
    Como $\snorm f = 0$, segue que $\pnorm p {g_\epsilon} < \epsilon$ e 
    $\pnorm r {h_\epsilon} < \epsilon$, ou seja, existem sequências 
    $(g_n)$ e $(h_n)$ com $f = g_n + h_n$ 
    para todo $n \in \bN$, tais que, para todo $\epsilon>0$ existe $N \in \bN$
    de modo que $\pnorm p {g_n}< \epsilon$ e $\pnorm r {h_n}<\epsilon$ 
    sempre que $n \ge N$. Desse modo, segue diretamente que $g_n$ e $h_n$
    convergem para $0$ nos seus respectivos espaços, e, já que $f = g_n + h_n$
    para todo $n \in \bN$, temos das propriedades de convergência em medida 
    que $f = 0$ $\mu$-q.t.p., portanto $[f] = [0]$.
\end{itemize}
\end{proof}

%ex 13 
\nextexe{quatro}
\textbf{Sejam $(X, \mcal F, \mu)$ e $f \in L^p(X, \mu) \cap L^\infty(X, \mu)$, 
para algum $1 \le p < +\infty$. Prove que 
\begin{align*}
  \pnorm \infty f = \lim_{q \rightarrow \infty} \pnorm q f.
\end{align*}}
\begin{proof}[\textbf{Dem:}] Fixado $p \in \bR$, 
  considere $q \in \bR$ tal que $1\le p < q < +\infty$.
  note que, $f \in L^p(X, \mu) \cap L^\infty(X, \mu) \subseteq L^q(X, \mu)$, e,
  já que $\mu(\set{x \in X \st |f(x)| > \pnorm \infty f}) = 0$, temos 
  $|f(x)| \le \pnorm \infty f$ $\mu$-q.t.p., desse modo,
  \begin{align*}
    \pnorm q f &= \cparth{\int |f|^q d\mu}^{1/q} = 
    \cparth{\int |f|^p|f|^{q-p}d\mu}^{1/q}
    \le \cparth{\int |f|^p \pnorm \infty f^{q-p} d\mu}^{1/q}\\
    &= \cparth{\pnorm \infty f^{q- p} \int |f|^p d\mu}^{1/q} 
    = \pnorm \infty f^{\frac{q-p}{q}} \pnorm p f^{p/q} 
    = \pnorm \infty f^{1 - p/q}\pnorm p f^{p/q}.
  \end{align*}
  Note que, para todo $k \in \bR$, $\lim\limits_{n \rightarrow \infty} k^{1/n} = 1$,
  de maneira análoga, segue que 
  \begin{align*}
    \lim\sup_{q}\pnorm q f 
    \le \lim\sup_{q} \pnorm \infty f^{1 - p/q} \pnorm p f^{p/q}
    =\lim_{q \rightarrow \infty} \pnorm \infty f^{1 - p/q} \pnorm p f^{p/q}
    = \pnorm \infty f. 
  \end{align*}
  Daqui, para cada $\delta > 0$ considere o conjunto 
  $E_{\delta} = \set{x \in X \st |f(x)| > \pnorm \infty f - \delta}$. 
  Da definição de $\pnorm \infty f$, temos que $\mu(E_{\delta})>0$. Com isso,
  dado $\delta > 0$,
  \begin{align*}
    \pnorm q f &= \cparth{\int \mparth f^q d\mu}^{1/q} \ge 
    \cparth{\int |f|^q \chi_{E_\delta}d\mu}^{1/q}
    \ge \cparth{\int_{E_\delta} \parth{\pnorm \infty f - \delta}^qd\mu}^{1/q}\\
    &= \cparth{\parth{\pnorm \infty f - \delta}^q \int_{E_\delta} d\mu}^{1/q}
    =\parth{\pnorm \infty f - \delta}\mu(E_\delta)^{1/q},
  \end{align*}
  já que $\lim_{q \rightarrow \infty} \mu(E_\delta)^{1/q} = 1$, temos 
  \begin{align*}
    \lim\inf_{q} \pnorm q f 
    &\ge \lim\inf_{q} \parth{\pnorm \infty f - \delta} \mu(E_\delta)^{1/q}
    = \parth{\pnorm \infty f - \delta}\lim_{q \rightarrow \infty} \mu(E_\delta)^{1/q}
    = \parth{\pnorm \infty f - \delta}.
  \end{align*}
  Como a desigualdade acima vale para todo $\delta > 0$, segue que 
  \begin{align*}
    \lim \sup_{q} \pnorm q f \le \pnorm \infty f \le \lim \inf_q \pnorm q f,
  \end{align*}
  ou seja,  $\lim \inf \pnorm q f = \lim \sup \pnorm q f 
  = \lim_{q\rightarrow \infty} \pnorm q f = \pnorm \infty f$.
\end{proof}

%ex 14 
\nextexe{quatro}
\textbf{Sejam $(X, \mcal F, \mu)$ e $ p \in (0, 1)$. Defina $L^p(X, \mu)$
de maneira semelhante ao caso $p \ge 1$ 
e mostre que a aplicação $d: L^p(X,\mu) \times L^p(X, \mu) \rightarrow \bR$
dada por 
\begin{align*}
  d(f, g) = \int \mparth{f - g}^p d\mu
\end{align*}
define uma métrica em $L^p(X, \mu)$ que é completa.}
\begin{proof}[\textbf{Dem:}] Primeiro, vejamos que 
  $d: L^p(X, \mu)\times L^p(X, \mu) \rightarrow \bR$ de fato 
  define uma métrica em $L^p(X, \mu)$. Dados 
  $f, g, h \in L^p(X, \mu)$,
\begin{itemize}
  \item $d(f, g) \ge 0$ e $d(f, g) = 0 \iff f = g.$\\
    Já que $|f-g|\ge 0$, segue diretamente que $\int \mparth{f - g}^p d\mu \ge 0$. 
    Agora , suponhamos $d(f, g) = 0$. Já que, 
    \begin{align*}
      d(f, g) = \int \mparth{f - g}^p d\mu = 0,
    \end{align*}
    e, $\mparth{f - g}^p$ é mensurável não-negativa, temos $\mparth{f-g}^p = 0$ 
    $\mu$-q.t.p. mais ainda, $f = g$ $\mu$-q.t.p, como queríamos. Suponhamos 
    agora que $f = g$. Então, 
    \begin{align*}
      d(f, g) = \int \mparth{f - g}^p d\mu = \int \mparth{f - f}^p d\mu = \int 0^p d\mu = 0.
    \end{align*}
  \item $d(f, g) = d(g, f)$. 
    Ora, 
    \begin{align*}
      d(f,g ) = \int \mparth{f - g}^p d\mu = \int \mparth{g - f}^p d\mu = d(g , f).
    \end{align*}
  \item $d(f, h) \le d(f, g) + d(g, h)$.\\
    Ora, se $f = h$, segue trivialmente a desigualdade. Suponhamos então 
    que $f \neq h$.
    Note que, para $p \in (0, 1)$ e $t \in (0, 1)$,
    \begin{align*}
      t^{1-p} \le 1^{1-p} = 1 \implies t \le t^p.
    \end{align*}
    Desse modo, 
    \begin{align*}
      t^p + (1-t)^p \ge t + 1 - t = 1.
    \end{align*}
    Com isso, vejamos que para $a, b \in \bR$ e $p \in (0, 1)$, 
    $|a + b|^p \le |a|^p + |b|^p$. Se $a = 0$ ou $b = 0$ ou $a + b = 0$, a desigualdade 
    segue de maneira trivial. Suponhamos que $a \neq 0$ e $b \neq 0$ e $a + b \neq 0$, 
    definindo $t = \frac{|a|}{|a| + |b|}$, temos $1 - t = \frac{|b|}{|a| + |b|}$. 
    Da desigualdade mostrada a pouco, 
    \begin{align*}
      &1 \le t^p + (1-t)^p = \parth{\frac{|a|}{|a|+|b|}}^p + \parth{\frac{|b|}{|a|+|b|}}^p
      \implies |a+b|^p\le ||a|+|b||^p \le |a|^p + |b|^p
    \end{align*}
    Desse modo, 
    \begin{align*}
      d(f, h) &= \int \mparth{f - h}^p d\mu = \int \mparth {f - g + (g - h)}^p d\mu  
      \le \int \mparth{f - g}^p + \mparth{g - h}^p d\mu\\ &= \int \mparth{f - g}^p d\mu 
      + \int \mparth{g - h}^p d\mu 
      = d(f, g) + d(g, h).
    \end{align*}
\end{itemize}
  Daqui, vejamos que o espaço $L^p(X, \mu)$ munido da métrica $d(\cdot, \cdot)$ 
  é um espaço completo. Dada uma sequência de Cauchy $(f_n)$, vejamos que 
  ela converge para algum elemento do espaço. Ora, para todo $\epsilon>0$, 
  existe $N \in \bN$ tal que, 
  \begin{align*}
    n, m \ge N \implies d(f_n, f_m) < \epsilon.
  \end{align*}
  Como visto em aula para o caso $p \ge 1$, para cada $k \in \bN$, 
  é possível criar subsequências $(g_k)$ de $(f_k)$ satisfazendo 
  \begin{align*}
    d(g_{k + 1}, g_k) < 2^{-k}, \forall k \in \bN
  \end{align*}
  Definindo $g: X \rightarrow \xbR$ de modo que 
  para todo $x \in X$, $g(x) = \mparth{g_1(x)} + \sum_{k \in \bN}\mparth{g_{k+1}(x)-g_k(x)}$,
  segue da construção que $g \in \mcal M^+(X, \mcal F)$. Pelo Lema de Fatou,
  \begin{align*}
    \int \mparth{g}^p d\mu &= \int \mparth{\lim_{l \rightarrow \infty} \mparth{g_1(x)} 
    + \sum_{k = 1}^{l}\mparth{g_{k+1}(x)-g_k(x)}}^pd\mu\\
    &\le \lim\inf_{l} \int \cparth{\mparth{g_1(x)}
    +\sum_{k=1}^l\mparth{g_{k+1}(x) - g_k(x)}}^pd\mu\\
    &\le \lim\inf_{l} \int \mparth{g_1(x)}^p d\mu + \sum_{k = 1}^l d(g_{k+1}, g_k)\\
    &\le \pnorm p {g_1}^p + \lim\inf_{l} \sum_{k=1}^l 2^{-k} = \pnorm p {g_1}^p + 1 \le \infty,
  \end{align*}
  ou seja, $g \in L^p(X, \mu)$, como queríamos.
\end{proof}


%ex 15
\nextexe{quatro}
\textbf{Sejam $(X, \mcal F, \mu)$ um espaço de medida, 
$(f_n)$ uma sequência em $L^p(X, \mu)$ que converge em quase todo ponto 
para uma função $f \in L^p(X, \mu)$, com $1\le p < +\infty$. Mostre que 
\begin{align*}
  \lim \pnorm p {f_n - f} = 0 \text{ se, e somente se, } \lim \pnorm p {f_n} = \pnorm p f.
\end{align*}}
\begin{proof}[\textbf{Dem:}].\\
  $(\implies)$ Suponhamos que $\lim \pnorm p {f_n - f} = 0$, então, 
  dado $\epsilon > 0$ existe $N \in \bN$ tal que para todo $n \ge N$, 
  \begin{align*}
    \pnorm p {f_n - f} < \epsilon.
  \end{align*}
  Da Desigualdade de Minkowski, temos $\pnorm p {f_n} \le \pnorm p {f_n - f} + \pnorm p f$
  e $\pnorm p f \le \pnorm p {f_n - f} + \pnorm p {f_n}$, desse modo, 
  \begin{align*}
    \mparth{\pnorm p {f_n} - \pnorm p f}\le \pnorm p {f_n - f} < \epsilon,
  \end{align*}
  isto é, $\lim \pnorm p {f_n} = \pnorm p f$.\\
  $(\impliedby)$ Suponhamos que $\lim \pnorm p {f_n} = \pnorm p f$, isto é,
  \begin{align*}
    \lim \int \mparth{f_n}^p = \int \mparth{f}^p d\mu.
  \end{align*}  
  Note que, $\mparth{f_n - f} \le \mparth{f_n} + \mparth{f} \le 2\max\set{\mparth{f_n}, \mparth{f}}$
  desse modo, 
  \begin{align*}
    \mparth{f_n - f}^p &\le \mparth{|f_n| + |f|}^p  
    \le 2^p \max\set{\mparth{f_n}, \mparth{f}}^p 
    \le 2^p (\mparth{f_n}^p + \mparth{f}^p).
  \end{align*}
  Defina para todo $n \in \bN$ a sequência $(g_n)$ satisfazendo 
  $g_n = 2^p(\mparth{f_n}^p + \mparth{f}^p) - \mparth{f_n - f}^p$, com isso, segue 
  por construção que $(g_n)$ é uma sequência em $\mcal M^+(X, \mcal F)$. Por hipótese, 
  sabemos que $(f_n)$ converge em quase todo ponto para $f \in L^p(X, \mu)$, sendo assim,
  $(g_n)$ converge para $2^{p+1}\mparth{f}^p$. Portanto, do Lema de Fatou, 
  \begin{align*}
    2^{p+1} \pnorm p f^p&= \int 2^{p+1}\mparth{f}^p d\mu =\int \lim \inf g_n d\mu \le
    \lim \inf \int g_n d\mu \\ 
    &\le \lim \inf \int 2^p(\mparth{f_n}^p + \mparth{f}^p) - \mparth{f_n - f}^p d\mu\\
    &= \lim \inf \cparth{2^p \parth{\int \mparth{f_n}^p d\mu + \int \mparth{f}^p d\mu}
    - \int \mparth{f_n - f}^p d\mu}\\
    &= 2^{p + 1} \pnorm p f^p - \lim\sup \pnorm p {f_n - f}^p. 
  \end{align*}
  Como $f \in L^p(X, \mu)$, 
  \begin{align*}
    2^{p+1}\pnorm p f^p \le 2^{p+1}\pnorm p f^p - \lim \sup \pnorm p {f_n - f}^p \implies
    \lim \sup \pnorm p {f_n - f}^p \le 0.
  \end{align*}
  Como $\pnorm p {f_n - f}^p \ge 0$ para todo $n \in \bN$, segue que 
  $\lim \sup \pnorm p {f_n - f}^p = 0 \ge \lim \inf \pnorm p {f_n - f}^p \ge 0$, 
  ou seja, $\lim \pnorm p {f_n - f}^p = 0$, sendo assim, $\lim \pnorm p {f_n - f} = 0$.
\end{proof}

%ex 16
\nextexe{quatro}
\textbf{Seja $(\bR, \mcal M, m)$ o espaço de medida de Lebesgue. Fixado $1 \le r < +\infty$, 
mostre que a função $f: \bR \rightarrow \bR$ dada por 
\begin{align*}
  f(x) = \parth{\frac{1}{x\mparth{\log x}^2}}^{1/r} \chi_{(0, 1/2) \cup (2, +\infty)}(x)
\end{align*}
pertence a $L^p(\bR, m)$ se, e somente se, $r=p$.}
\begin{proof}[\textbf{Dem:}]
  Suponhamos que $f \in L^p(\bR, m)$, então, 
  \begin{align*}
    \int f dm &= \int \parth{\frac{1}{x\mparth{\log x}^2}}^{1/r} 
    \chi_{(0, 1/2) \cup (2, +\infty)}(x) dm \\
    &= \int_{(0, 1/2)} \parth{\frac{1}{x\mparth{\log x}^2}}^{1/r}dm
    + \int_{(2, \infty)} \parth{\frac{1}{x\mparth{\log x}^2}}^{1/r} dm
  \end{align*}
\end{proof}

%ex 17 
\nextexe{quatro}
\textbf{Sejam $(X, \mcal F ,\mu)$ um espaço de medida e 
$1 \le r, p, q \le +\infty$ tais que $\frac{1}{p} + \frac{1}{q} = \frac{1}{r}$. 
Prove que se $f \in L^p(X, \mu)$ e $g \in L^q(X, \mu)$, então $fg \in L^r(X, \mu)$ 
\begin{align*}
  \pnorm r {fg} \le \pnorm p f \pnorm q g.
\end{align*}
Generalize para uma soma finita $\sum \frac{1}{p_k} = \frac{1}{r}$.}
\begin{proof}[\textbf{Dem:}] Note que, se $p = q = \infty$ temos $r = \infty$,
  e por conseguinte, a 
  desigualdade segue de maneira direta. Mais ainda, se $q = \infty$ e $p \neq \infty$,
  temos $p = r \neq \infty$. Por definição, $g \le \pnorm \infty g$ $mu$-q.t.p., 
  portanto, 
  \begin{align*}
    \int \mparth{fg}^r d\mu \le \pnorm \infty g^r \int \mparth{f}^r d\mu, 
  \end{align*}
  Como $p = r$,
  \begin{align*}
    \pnorm r {fg}= \parth{\int \mparth{fg}^r d\mu}^{1/r} 
    \le \pnorm \infty g \parth{\int \mparth{f}^r d\mu}^{1/r} = \pnorm \infty g \pnorm r f 
    = \pnorm \infty g \pnorm p f.
  \end{align*}
  Sendo assim, suponhamos $p \neq \infty$ e $q \neq \infty$. Note que se $f= 0$ ou $g =0$
  a desigualdade segue diretamente. Com isso, suponhamos também que $f\neq 0$ e $g \neq 0$. 
  Daqui, visto que 
  $\frac{r}{p} + \frac{r}{q} = 1$, ou seja, $u = \frac{p}{r}$ é conjugado de 
  $v = \frac{q}{r}$, da Desigualdade de Young, definindo $a = \frac{|f|^r}{\pnorm p f^r}$
  e $b =\frac{|g|^r}{\pnorm q g^r}$,
  \begin{align*}
    \int \frac{\mparth{fg}^r}{\pnorm p f^r \pnorm q g^r} d\mu &\le 
    \int \frac{1}{u}\parth{\frac{\mparth{f}^r}{\pnorm p f^r}}^u +
    \frac{1}{v}\parth{\frac{\mparth{g}^r}{\pnorm q g^r}}^v d\mu \\
    &= \frac{r}{p} \frac{\int \mparth{f}^p d\mu}{\pnorm p f^p}  
    + \frac{r}{q} \frac{\int \mparth{g}^q d\mu}{\pnorm q g^q} \\ 
    &= \frac{r}{p} + \frac{r}{q} = 1. 
  \end{align*}
  Ou seja, 
  \begin{align*}
    \int \frac{\mparth{fg}^r}{\pnorm p f^r \pnorm q g^r} d\mu \le 1 
    &\implies \int \mparth{fg}^r d\mu \le \pnorm p f^r \pnorm q g^r \\
    &\implies \pnorm r {fg} \le \pnorm p f \pnorm q g. 
  \end{align*}
  Para a generalização utilizando a soma finita 
  $\sum\limits_{k = 1}^n \frac{1}{p_k} = \frac{1}{r}$, com $1\le p_k \le +\infty$
  para todo $k \in [1, n] \cap \mbb Z$,
  seguem as desigualdades para o caso $n \le 2$. Suponhamos por indução em $n$ que, dada 
  uma sequência $(f_k)_{k=1}^n$ de funções satisfazendo $f_k \in L^{p_k}(X, \mu)$ para 
  todo $k \in [1, n] \cap \mbb Z$, 
  para todo inteiro $m$ satisfazendo $1 \le m < n$, temos 
  $\pnorm r {\prod\limits_{k = 1}^m f_k} \le \prod\limits_{k = 1}^m \pnorm {p_k} {f_k}$
  e $\prod_{k = 1}^m f_k \in L^r(X, \mu)$ sempre
  que $\sum\limits_{k = 1}^{m}\frac{1}{p_k} = \frac{1}{r}$. Desse modo, 
  se $p_n = \infty$, 
  \begin{align*}
    \int \mparth{\prod_{k = 1}^n f_k}^r d\mu 
    \le \pnorm \infty {f_n}^r \int \mparth{\prod_{k = 1}^{n-1} f_k}^r d\mu
    \le \pnorm \infty {f_n}^r \parth{\prod_{k = 1}^{n - 1}\pnorm {p_k}{f_k}}^r = 
    \parth{\prod_{k=1}^n \pnorm {p_k}{f_k}}^r.
  \end{align*}
  Agora, suponha $p_n \neq \infty$. Definindo 
  $\frac{1}{v_1} = \sum\limits_{k = 1}^{n-1} \frac{1}{p_k}$
  e $\frac{1}{v_2} = \frac{1}{p_n}$, temos $\frac{1}{v_1}+ \frac{1}{v_2} = \frac{1}{r}$. 
  Mais ainda, defina $f = f_n$ e $g = \prod\limits_{k=1}^{n-1} f_k$ e note que 
  $|f|$ e $|g|$ são mensuráveis não-negativos.  
  Como $\sum\limits_{k=1}^{n-1} \frac{1}{p_k} = \frac{1}{v_1}$ e 
  $f_k \in L^{p_k}(X, \mu)$ para $k \in [1, n]\cap \mbb Z$, da hipótese de indução 
  segue que $\prod\limits_{k=1}^{n-1} f_k \in L^{v_1}(X, \mu)$, e,
  \begin{align*}
    \pnorm {v_1} g = \pnorm {v_1} {\prod_{k=1}^{n-1} f_k} 
    \le \prod_{k=1}^{n-1}\pnorm {p_k} {f_k} 
  \end{align*}  
  Com isso, da primeira parte do 
  exercício, temos $\prod_{k=1}^n f_k = fg \in L^r(X, \mu)$, e,
  \begin{align*}
    \pnorm r {\prod_{k = 1}^n f_k} = \pnorm r {fg} \le 
    \pnorm {v_2} f \pnorm {v_1} g 
    \le \pnorm {p_n} {f_n} \prod\limits_{k=1}^{n-1} \pnorm {p_k}{f_k} = 
    \prod_{k = 1}^n \pnorm {p_k} {f_k},
  \end{align*}
  como queríamos.
\end{proof}

%ex 18
\nextexe{quatro}
\textbf{Sejam $(X, \mcal F, \mu)$ um espaço de medida e $1 \le r, p, q \le +\infty$
tais que $\frac{1}{p} + \frac 1 q + \frac 1 r = 1$. Prove que 
se $f \in L^p (X, \mu)$, $g \in L^q(X, \mu)$ e $h \in L^r(X, \mu)$, 
então $fgh \in L^1(X, \mu)$ e 
\begin{align*}
  \pnorm 1 {fgh} \le \pnorm p f \pnorm q g \pnorm r h.
\end{align*}
Generalize para uma soma finita $\sum \frac{1}{p_k} = 1$.}
\begin{proof}[\textbf{Dem:}]Note que o resultado segue de maneira direta pelo 
  exercício anterior. 
\end{proof}

%ex 19
\nextexe{quatro}
\textbf{Sejam $(X, \mcal F, \mu)$ um espaço de medida, considere 
$p \in \bR$ tal que $1 \le p < +\infty$ e $q$ o seu conjugado. Prove que para 
toda $f \in L^p(X, \mu)$ vale que 
\begin{align*}
  \pnorm p f = \sup \set{\int \mparth{fg} d\mu \st g\in L^q(X, \mu) \text{ e } 
  \pnorm q g \le 1}.
\end{align*}}
\begin{proof}[\textbf{Dem:}] Ora, dados 
  $f \in L^p(X, \mu)$ e qualquer $g \in L^q(X, \mu)$ satisfazendo $\pnorm q g \le 1$,
  temos, 
  \begin{align*}
    \pnorm 1 {fg} \le \pnorm p f \pnorm q g \le \pnorm p f.
  \end{align*}
  Desse modo, 
  \begin{align*}
    \sup\set{\pnorm 1 {fg} \st g\in L^q(X, \mu), \pnorm q g \le 1} \le \pnorm p f.
  \end{align*}
  Mais ainda, note que a função conjugada de $f$ definida como 
  \begin{align*}
    f^* \coloneq \pnorm p f^{1-p} \sgn(f)\mparth{f}^{p-1},
  \end{align*}
  em que $\sgn(f)(x) = \begin{cases}
    \frac{f(x)}{\mparth{f(x)}}, & \text{se } f(x) \neq 0;\\
    0, &\text{se } f(x) = 0.
  \end{cases}$, satisfaz as seguintes propriedades: $f^* \in L^q(X, \mu)$, 
  $\int f\cdot f^* d\mu = \pnorm p f$ e $\pnorm q {f^*} = 1$. Ou seja, 
  \begin{align*}
    \pnorm p f = \int f\cdot f^* d\mu \le \sup \set{\int \mparth{fg}d\mu \st 
    g \in L^q(X, \mu), \pnorm q g \le 1}.
  \end{align*}
  Com isso, temos $\pnorm p f = \sup\set{\int \mparth{fg} d\mu st g \in L^q(X, \mu), 
  \pnorm q g \le 1}$.
\end{proof}

%ex 20 
\nextexe{quatro}
\textbf{Sejam $(X, \mcal F, \mu)$ um espaço de medida e $1 < p < +\infty$. 
Mostre que o espaço $L^p(X, \mu)$ é estritamente convexo, ou seja, 
se $f, g \in L^p(X, \mu)$ são tais que $\pnorm p f = 1$, $\pnorm p g = 1$ e 
$f \neq g$, então 
\begin{align*}
  \pnorm p {(1-\lambda)f + \lambda g}< 1, \quad \forall \lambda \in (0, 1).
\end{align*}
Mostre também, através de contraexemplos, que o resultado não é válido 
para $p=1$ e $p = +\infty$ no espaço de medida de Lebesgue em $\bR$.}
\begin{proof}[\textbf{Dem:}] Visto que 
  $f, g \in L^p(X, \mu)$, então 
  $(1 - \lambda)f, \lambda g, (1-\lambda)f + \lambda g \in L^p(X, \mu)$.
  Desse modo, da desigualdade de Minkowski, para todo $\lambda \in (0, 1)$, temos
  \begin{align*}
    \pnorm p {(1-\lambda)f + \lambda g} &\le \pnorm p {(1-\lambda) f} + \pnorm p {\lambda g}
    = \mparth{1 - \lambda} \pnorm p f + \mparth{\lambda} \pnorm p g
    = \mparth{1 - \lambda} + \mparth{\lambda} = 1. 
  \end{align*}
  Vejamos que $\pnorm p {(1 - \lambda) f + \lambda g} \neq 1$. Note que a igualdade
  é satisfeita se, e somente se, $|f| = \alpha |g|$ $\mu$-q.t.p. para algum  $\alpha \ge 0$. 
  Entretanto, 
  \begin{align*}
    1 = \pnorm p f = \alpha \pnorm p g = \alpha,
  \end{align*}
  isto é, $f = g$, o que contradiz o fato de que $f \neq g$.\\
  Considere o seguinte contraexemplo para $p=1$: Sejam $n \in \bN$, 
  $f = n \chi_{(1/n, 2/n)}$, $g = n \chi_{(2/n, 3/n)}$ 
  e $\lambda \in (0 , 1)$, então
  \begin{align*}
    \int \mparth{(1- \lambda) f + \lambda g}^p dm 
    &= \int (1-\lambda) n \chi_{(1/n, 2/n)} + \lambda n\chi_{(2/n, 3/n)} dm\\ 
    &= (1-\lambda) n \cdot m((1/n, 2/n)) + \lambda n \cdot m((2/n, 3/n))\\ 
    &= (1 - \lambda) + \lambda = 1
  \end{align*}
  Considere o seguinte contraexemplo para $p = \infty$: Seja $f=\chi_{[2, 4)}$
  e $g = \chi_{(1, 3]}$. Note que, 
  \begin{align*}
    \pnorm \infty f = \inf \set{\alpha\ge 0 \st 
    \m\parth{\set{x \in X \st \mparth{f(x)} > \alpha}} = 0} = \inf A = 1,
  \end{align*}
  Pois, para todo $x \ge 1$, temos $x \in A$ e $x < 1$, segue que $x \not \in A$. 
  Semelhantemente, segue que $\pnorm \infty g = 1$. Com isso, 
  \begin{align*}
    \pnorm \infty {(1-\lambda) \chi_{(1, 3]}+ \lambda \chi_{[2, 4)}} 
    = \pnorm \infty { (1-\lambda)\chi_{(1, 2)} + \chi_{[2, 3]} + \lambda \chi_{(3, 4)}} = 1.
  \end{align*}
\end{proof}

%ex 21 
\nextexe{quatro}
\textbf{Sejam $(\bR, \mcal M, m)$ o espaço de medida de Lebesgue e $f: \bR \rightarrow \bR$
uma função mensurável tal que $f \not \in L^\infty(\bR, m)$. 
\begin{enumerate}[label=\alph*)]
  \item Defina os conjuntos de nível $A_n = \set{ x \in \bR \st \mparth{f(x)} > n}$. 
    Mostre que $m(A_n) > 0$, para todo $n \in \bN$.
    \begin{proof}[\textbf{Dem:}]
      Como $f \not \in L^\infty(\bR, m)$, segue que para todo $\alpha \ge 0$, 
      $m(\set{x \in \bR \st \mparth{f(x)} > \alpha}) > 0$, com isso, 
      segue diretamente que para todo $n \in \bN$, $m(A_n) > 0$.
    \end{proof}
  \item Considere os conjuntos $D_n = \set{x \in \bR \st n < \mparth{f(x)} \le n+1}$. 
    Mostre que existe uma sequência estritamente crescente de índices $(n_k)_{k \in \bN}$
    tal que $m(D_{n_k}) > 0$ para todo $k$.
    \begin{proof}[\textbf{Dem:}] Note que, 
      para todo $k \in \bN$, temos $A_k = \bigcup_{j = k}^\infty D_j$. Como 
      $m(A_k)> 0, \forall k \in \bN$, existe $j \in [k, \infty)\cap \bN$ 
      tal que $m(D_j) > 0$. Para 
      cada $k \in \bN$, defina $n_1 = \min \set{j \in [1, \infty)\cap \bN \st m(D_j) > 0}$ 
      e, para $k \neq 1$, $n_k = \min \set{j \in (n_{k-1}, \infty)\cap \bN \st m(D_j) > 0}$.
      Note que cada $n_k$ está bem definido, pois, 
      $A_{n_{k}+1} = \bigcup_{j = n_{k}+1}^\infty D_j$ e 
      $m(A_{n_{k}+1}) > 0$, portanto, existe $j \in (n_{k}, \infty) \cap \bN$ 
      tal que $m(D_j)> 0$.
      Com isso, segue por construção que para todo $k \in \bN$, temos $n_k < n_{k+1}$ e 
      $m(D_{n_k}) > 0$.
    \end{proof}
  \item Mostre que é possível escolher subconjuntos mensuráveis $F_k \subset D_{n_k}$
    tais que $0 < m(F_k) < \frac{1}{2^k}$.
    \begin{proof}[\textbf{Dem:}]
     aa 
    \end{proof}
  \item Defina a sequência $E_k = \bigcup_{j = k}^\infty F_j$. Prove 
    que $\set{E_k}_{k \in \bN}$ satisfaz simultaneamente as seguintes 
    condições: 
    \begin{enumerate}[label=\roman*)]
      \item $0 < m(E_k) < +\infty$;
      \item $E_{k+1} \subset E_k$;
      \item $m(E_k \setminus E_{k+1}) > 0$;
      \item $\lim m(E_k) = 0$;
      \item $E_k \subset \set{x \in \bR \st \mparth{f(x)} > k}$.
    \end{enumerate}
\end{enumerate}}
